\chapter{Logística da Distribuição de Vacinas}

\section{Monitoramento de Estoques}
Além do transporte, a gestão de estoques desempenha um papel crucial para evitar desperdícios. No Brasil, sistemas informatizados como o Sistema de Informações do Programa Nacional de Imunizações (SIPNI) são utilizados para monitorar estoques em tempo real, permitindo que os gestores ajustem a distribuição de doses conforme a demanda de cada região \cite{frazzon:2024}.

\section{Etapas da Distribuição de Vacinas}
No processo de distribuição de vacinas, o primeiro passo é a distribuição de vacinas pelo Ministério da Saúde a partir do centro do Centro Nacional de Distribuição de Insumos (CENADI), que é localizado em  São Paulo/SP. Nesse centro são coletados amostras de todos os lotes de vacina, que são enviados para Instituto Nacional de Controle de Qualidade em Saúde (INCQS), para que seja 
comprovada a qualidade dos produtos.

O passo seguinte é a distribuição da quantidade de doses definida proporcionalmente para cada estado e esta operação é encerrada em até 48 horas. Após a chegada, geralmente nas capitais, é responsabilidade dos estados seguir com a distribuição para os municípios \cite{comprovei:2024}.

Cabe ressaltar que, dependendo do estado, pode haver uma etapa intermerdiária entre o depósito estadual e os municípios. Estados de grande dimensão geográfica, como o de Minas Gerais, ainda realizam a etapa de distribuição do depósito estadual para Superintendências Regionais de Saúde (SRS) \cite{LimaEtAl:2019}). Por outro lado, estados de malha compacta e centralizada, como o Rio de Janeiro, operam majoritariamente por meio de fluxos diretos do depósito estadual para os municípios.

A distribuição aos municípios ocorre a partir de uma Central Municipal de Rede de Frio e deve ser realizada em até 7 dias e cada município define a estratégia de distribuição nos postos de saúde \cite{canalgov:2021}.

Por fim, os municípios relataram que em média são gastos 30 minutos, desde a chegada do caminhão, conferência das vacinas e armazenagem na câmara fria \cite{LimaEtAl:2019}.

\section{Armazenagem}
Para manter a eficácia dos imunizantes, as vacinas devem ser armazenadas em temperaturas específicas, que variam de acordo com o seu tipo. A maioria deve ficar entre 2°C e 8°C, mas algumas precisam ser mantidas em temperaturas ainda mais baixas, como -15°C ou -70°C. Por isso, durante o transporte, as vacinas devem ser mantidas em uma embalagem térmica apropriada. A embalagem escolhida precisa ser validada antes do transporte para garantir sua eficácia e a temperatura também deve ser verificada regularmente durante o trajeto \cite{bosch:2024}.

A tecnologia de monitoramento, com sensores térmicos e dataloggers instalados nas maletas térmicas e nos caminhões, permite o controle preciso da temperatura em tempo real, garantindo que as condições ideais sejam mantidas durante todo o trajeto \cite{frazzon:2024}. Essas informações são analisadas pela Anvisa (Agência Nacional de Vigilância Sanitária) para a posterior liberação dos imunizantes \cite{comprovei:2024}.

\section{Modais de transporte}
No Brasil, para realizar o transporte de produtos farmacêuticos, é necessário que a transportadora obtenha a Autorização de Funcionamento de Empresa (AFE) junto à Anvisa, podendo ainda ser necessário uma Autorização Especial (AE) caso haja substâncias especiais.

O envio de imunizantes para os Estados é realizado por transporte rodoviário ou aéreo.

O transporte rodoviário de vacinas exige veículos equipados com sistemas de refrigeração para manter a temperatura adequada. Esses veículos são monitorados continuamente durante o trajeto para garantir que a temperatura permaneça dentro dos limites especificados.

O transporte aéreo é exige aeronaves que dispõem de controle de temperatura no porão, capazes de atender as diferentes faixas de temperatura que medicamentos e vacinas requerem, como de 2 a 8°C, 15 a 25°C ou 2 a 25°C. \cite{aglcargo:2025}
\bigskip
